\anexo{Cronograma de Actividades}

Las columnas del cronograma representan los cuartos del calendario planificado para todo el PFI (25\,\%, 50\,\%, 75\,\% y 100\,\% del tiempo total del proyecto), y el sombreado indica en qué cuarto se trabajó cada actividad, en base al historial real de desarrollo (commits y bitácoras de sesión). Los cuartos ya transcurridos (25\,\%, 50\,\% y 75\,\%) tienen fecha calendario asociada; el cuarto 100\,\% todavía no tiene fecha de entrega fijada por la cátedra al momento de esta entrega, y se deja indicado como ``a definir'' en lugar de estimarse arbitrariamente. Esta entrega corresponde al hito del 75\,\%: además de las actividades ya en curso desde el cuarto anterior, se sumaron el entrenamiento y evaluación del primer modelo de Machine Learning supervisado, el desarrollo del modelo de negocio, la viabilidad económica y el branding, y -- a partir de observaciones puntuales de la revisión oral del hito del 50\,\% -- un sistema de cuentas y autenticación, la resolución de la identidad de producto por catálogo, la agregación del score entre vendedores y un endurecimiento de seguridad sobre las cuentas. El análisis final de resultados y la redacción completa del informe (incluyendo Resumen, Abstract y Agradecimientos) corresponden a la última etapa del proyecto y todavía no se iniciaron.

\newcommand{\avance}{\cellcolor{black!75}}

\begin{table}[ht]
	\centering
	\small
	\renewcommand{\arraystretch}{0.85}
	\caption{Cronograma de actividades}
	\label{tab:gantt}
	\begin{tabularx}{\textwidth}{@{}X c c c c@{}}
		\toprule
		\textbf{Actividad} & \textbf{\footnotesize 25\% (Abr--May 26)} & \textbf{\footnotesize 50\% (Jun--Ago 26)} & \textbf{\footnotesize 75\% (Sep 26)} & \textbf{\footnotesize 100\% (a definir)} \\
		\midrule
		Definición del problema y alcance                    & \avance &         &         &         \\
		Marco teórico y antecedentes académicos               & \avance & \avance &         &         \\
		Estado del arte y comparación de herramientas          & \avance & \avance &         &         \\
		User Research y validación inicial                    & \avance & \avance &         &         \\
		Definición de arquitectura general                    & \avance & \avance &         &         \\
		Diseño del modelo de datos en PostgreSQL               & \avance & \avance &         &         \\
		Conexión con API oficial de Mercado Libre              & \avance & \avance &         &         \\
		Construcción del pipeline de captura de datos          & \avance & \avance &         &         \\
		Construcción del dataset histórico propio              & \avance & \avance & \avance &         \\
		Desarrollo del backend con Next.js (API Routes)        & \avance & \avance &         &         \\
		Implementación del motor analítico                    &         & \avance &         &         \\
		Desarrollo del dashboard web                          & \avance & \avance &         &         \\
		Cálculo de métricas temporales                        &         & \avance &         &         \\
		Entrenamiento y evaluación del modelo de ML supervisado &         &         & \avance &         \\
		Modelo de negocio, viabilidad económica y branding     &         &         & \avance &         \\
		Sistema de cuentas, autenticación y alertas            &         &         & \avance &         \\
		Identidad de producto por catálogo y score agregado    &         &         & \avance &         \\
		Endurecimiento de seguridad (rate limiting, cambio de contraseña) &  &        & \avance &         \\
		Pruebas funcionales y ajustes                         & \avance & \avance & \avance &         \\
		Análisis de resultados y limitaciones                 &         &         & \avance &         \\
		Redacción final y preparación de defensa               &         &         &         &         \\
		\bottomrule
	\end{tabularx}
\end{table}